\section*{Filtro Bandpass (retroazione multipla)}

\noindent\colorbox{orange}{
  \begin{minipage}{1.0\linewidth}
    \paragraph{Richiesta (a)}
    Dimostrare, sulla base di semplici argomenti di natura qualitativa, che il circuito proposto
    svolga effettivamente la funzione di filtro passa-banda, discutendone la risposta ai limiti
    del dominio delle frequenze reali.
  \end{minipage}
}
\begin{itemize}
\item
  Per $f \to 0$ i condensatori $C_1$ e $C_2$ aprono i loro rami, le tensioni d'ingresso all'op-amp
  sono nulle
  \begin{equation*}
    V_+ = 0, ~V_- = 0 \implies \colorbox{blu}{\boxed{V_{OUT} = 0}}
  \end{equation*}
\item
  Per $f \to \infty$ il condensatore $C_2$ mette in ``corto-circuito'' $V_A$ con $V_{OUT}$
  \begin{equation*}
    V_A = V_{OUT}
  \end{equation*}
  il condensatore $C_1$ mette in ``corto-circuito'' $V_A$ con $V_-$
  \begin{equation*}
    V_{OUT} = A_d(0-V_A) = A_d(0-V_{OUT}) \implies \colorbox{blu}{\boxed{V_{OUT} = \frac{0}{1 + A_d} = 0}}
  \end{equation*}
\end{itemize}

\noindent\colorbox{orange}{
  \begin{minipage}{1.0\linewidth}
    \paragraph{Richiesta (b)}
    Derivare analiticamente la funzione di trasferimento complessa $A_v(s)=V_{OUT}(s)/V_S(s)$
    nell’ipotesi che l’operazionale sia ideale (\textit{suggerimento: si determini la tensione nel
      punto A del circuito in termini di $V_S$ e $V_{OUT}$; quindi si esprima $V_A$ in termini di $V_{OUT}$
      assumendo che l'op-amp funzioni in regime lineare ed utilizzando il principio di cortocircuito
      virtuale; si eguaglino infine le espressioni ottenute per ottenere la funzione di trasferimento
    }).
    Per semplicità di calcolo, si assuma inoltre (ipotesi compatibile con i
    valori e le incertezze nominali dei componenti) che $R_1=R_2=R$, $R_3=100R$, $C_1=C_2=C$.
  \end{minipage}
}

\paragraph{Risposta (b)}
Seguendo il suggerimento della richiesta, attraverso la \textit{KCL}
si ottiene la seguente relazione per la tensione $V_A$ in termini di $V_S$
e $V_{OUT}$
\begin{equation*}
  \frac{V_S - V_A}{R} = sC(V_A - V_{OUT}) + sC(V_A - V_-) + \frac{V_A}{R}
\end{equation*}
attraverso la \textit{KCL} si ottiene anche una seconda relazione per $V_A$
in termini di $V_{OUT}$
\begin{equation*}
  sC(V_A - V_-) = \frac{V_- - V_{OUT}}{100R}
\end{equation*}
siccome $V_+ = 0$ nell'ipotesi di validit\'a per il principio di corto-cicuito
virtuale vale
\begin{equation*}
  V_+ = V_- \implies sCV_A = -\frac{V_{OUT}}{100R}
\end{equation*}
la soluzione del sistema dell'equazione precedente appaiata con la prima espressione
ricavata coinvolge \textit{unicamente} $V_S$ e $V_{OUT}$. Il primo passaggio
\'e una manipolazione algebrica della prima espressione e la sostituzione di $V_- = 0$
\begin{equation*}
  V_S - V_A = sRC(V_A - V_{OUT}) + sRCV_A + V_A
\end{equation*}
si raccolgono a destra i termini con $V_A$
\begin{equation*}
  V_S = 2(1 + sRC)V_A - sRCV_{OUT}
\end{equation*}
si sostituisce $V_A$ e raccoglie nuovamente
\begin{equation*}
  V_S = -\frac{1}{50sRC}\left(1 + sRC + 50s^2R^2C^2\right)V_{OUT}
\end{equation*}
finalmente si ottiene la funzione di trasferimento
\begin{equation}
  \label{eq:bp-multi-feedback-tfunc}
  \colorbox{blu}{\boxed{H(s) = -\frac{50sRC}{1 + sRC + 50s^2R^2C^2}}}
\end{equation}

\noindent\colorbox{orange}{
  \begin{minipage}{1.0\linewidth}
    \paragraph{Richiesta (c)}
    Dalla funzione di trasferimento dedurre i valori attesi per i parametri del filtro (frequenza
    centrale, guadagno massimo e larghezza banda) e confrontarli con le misure di cui ai
    punti 1b, 1c, 1d.
  \end{minipage}
}
\paragraph{Risposta (c)}
La funzione di trasferimento \ref{eq:bp-multi-feedback-tfunc} \'e in forma canonica, calcolando
le parti reali ed immaginarie del numeratore e denominatore con dominio ristretto all'asse
immaginario si ottiene
\begin{equation*}
  \begin{cases}
    \mathcal{R}_{\omega}^N = 0 \\
    \mathcal{I}_{\omega}^N = -50RC\omega \\
    \mathcal{R}_{\omega}^D = 1 - 50R^2C^2\omega^2 \\
    \mathcal{I}_{\omega}^D = RC\omega
  \end{cases}
\end{equation*}
di conseguenza ottengo la funzione di guadagno mediante \ref{eq:da-tfunc-a-gain}
\begin{equation*}
  G(\omega) = \frac{50RC\omega}{\sqrt{(1 - 50R^2C^2\omega^2)^2 + (RC\omega)^2}}
\end{equation*}
dividendo numeratore e denominatore per $RC\omega$ si ottiene un espressione pi\'u
semplice da cui ricavare il guadagno massimo
\begin{equation*}
  G(\omega) = \frac{50}{\sqrt{1 + \left(\frac{1 - 50R^2C^2\omega^2}{RC\omega}\right)^2}}
\end{equation*}
Il valore minimo che l'espressione sotto la radice potrebbe idealmente assumere \'e $1$,
questo valore \'e ottenibile per $1 - 50R^2C^2\omega^2 = 0$ dunque
\begin{equation}
  \label{eq:bp-multi-feedback-gain-max}
  \colorbox{blu}{\boxed{\max_{\omega \in [0, \infty)}G(\omega) = \frac{50}{\sqrt{1 + 0}} = 50}}
\end{equation}
inoltre attraverso la sostituzione $\omega \to f = 2\pi \omega$ si ottiene la frequenza centrale
\begin{equation}
  \label{eq:bp-multi-feedback-frequenza-centrale}
  \colorbox{blu}{\boxed{f = \frac{1}{10\pi\sqrt{2}RC} \approx 2250~\text{Hz}}}
\end{equation}

\noindent\colorbox{orange}{
  \begin{minipage}{1.0\linewidth}
    \paragraph{Richiesta (d)}
    Dare un’interpretazione qualitativa (e se possibile quantitativa) delle caratteristiche del
    segnale di uscita osservato al punto 1e. Che tipo di segnale appare? A quale ampiezza?
  \end{minipage}
}

\paragraph{Risposta (d)}
Se si invia nell'ingresso un onda quadra di frequenza pari alla frequenza di centro-banda, siccome
l'onda quadra \'e rappresentabile come una combinazione lineare di onde sinusoidali la cui prima
componente a frequenza non-nulla \'e a frequenza di centro-banda. Il risultato dovrebbe essere
un unica onda sinusoidale
\begin{equation}
  \colorbox{blu}{\boxed{V_{OUT} = 50\cdot\frac{4\cdot 50~\text{mV}}{\pi} \approx 3.18~\text{V}}}
\end{equation}

